%%%%%%%%%%  scan idx 371  |  folio 358  %%%%%%%%%%
ce q\hphantom{u}i
montre que \(u^{-1}\circ\Gr(\widetilde c)=\Id\).

\textbf{Corollaire 6.12.} Si \(K\) est de caractéristique nulle, et \(\Lambda\) une \(K\)-\(\lambda\)-algèbre de filtration discrète, alors la classe de Chern complétée :
\[
  \widetilde c:\Lambda\longrightarrow \Ch(\Gr\Lambda)
\]
est un isomorphisme.

La démonstration est contenue dans celle de 6.11.

\textbf{Définition 6.13.} Soit \(A\) une \(K\)-algèbre graduée, et \(\Lambda\) une \(K\)-\(\lambda\)-algèbre augmentée~; on appelle \emph{théorie de Chern de \(\Lambda\) à valeurs dans \(A\)} tout \(\lambda\)-homomorphisme : \(\Lambda\longrightarrow \Ch(A)\), compatible avec les augmentations.

On définit un foncteur sur la catégorie des \(K\)-algèbres graduées, à valeurs dans la catégorie des ensembles, en associant à toute \(K\)-algèbre graduée \(A\) l'ensemble des théories de Chern de \(\Lambda\) à valeurs dans \(A\), soit \(\Hom_\lambda(\Lambda,\Ch(A))\).

\textbf{Proposition 6.14.} Le foncteur défini en 6.13 est représentable. Si \(K\) est de caractéristique nulle, et \(\Lambda\) de filtration discrète, il est représenté par le gradué associé à \(\Lambda\).

Montrons d'abord le premier point. On considère la \(K\)-algèbre graduée engendrée par des générateurs \(c^i(x)\), avec \(i\ge1\), \(x\in\Fil^1\Lambda\), \(c^i(x)\) étant de degré \(i\), et on fait son quotient par les relations suivantes :
\[
\begin{array}{ll}
\text{i)}&
 c^k(x+y)=\displaystyle\sum_{i+j=k}c^i(x)c^j(y),
 \quad\text{pour } k\ge1,\ x,y\in\Fil^1\Lambda\ ;\\[0.8em]
\text{ii)}&
 c^k(xy)=P''_k(c^i(x),c^j(y))
 \quad(\text{cf. }6.1),
 \quad\text{pour } k\ge1,\ x,y\in\Fil^1\Lambda\ ;\\[0.4em]
\text{iii)}&
 c^k(ax)=R_{k,a}(c^i(x)
 \quad(\text{cf. }6.1),
 \quad\text{pour } k\ge1,\ a\in K,\ x\in\Fil^1\Lambda\ ;\\[0.4em]
\text{iv)}&
 c^k(\gamma^j(x))=Q''_{j,k}(c^i(x))
 \quad(\text{cf. }6.1),
 \quad\text{pour } j,k\ge1,\ x\in\Fil^1\Lambda,
\end{array}
\]
